\documentclass[11pt,a4paper]{ctexart}
\usepackage[margin=21mm,headheight=23pt]{geometry}
\usepackage{xcolor,booktabs,tabularx,enumitem,fancyhdr,hyperref,graphicx}
\definecolor{ink}{HTML}{173B2B}
\definecolor{accent}{HTML}{5E874A}
\definecolor{muted}{HTML}{5F7167}
\definecolor{line}{HTML}{DCE6D7}
\hypersetup{colorlinks=true,linkcolor=ink,urlcolor=accent}
\setlength{\parindent}{0pt}
\setlength{\parskip}{6pt}
\setlist[itemize]{leftmargin=*,itemsep=3pt,topsep=3pt}
\pagestyle{fancy}
\fancyhf{}
\fancyhead[L]{\small\color{muted}REINS / 客户内测}
\fancyhead[R]{\small\color{muted}2026 年 10 月}
\fancyfoot[R]{\small\color{muted}\thepage}
\renewcommand{\headrulewidth}{0.4pt}
\renewcommand{\headrule}{\hbox to\headwidth{\color{line}\leaders\hrule height \headrulewidth\hfill}}
\newcommand{\sectiontitle}[1]{\vspace{5pt}{\large\bfseries\color{ink}#1}\par\vspace{2pt}}
\begin{document}
\vspace*{4mm}
{\fontsize{31}{38}\selectfont\bfseries\color{ink}让每一次 Agent 任务\\都能算清、看明白。\par}
\vspace{5mm}
{\large\color{muted}Reins 客户试用手册：从一个已有 Python 工作流开始}\par
\vspace{6mm}
\colorbox{line}{\parbox{0.91\linewidth}{\vspace{5pt}\textbf{第一步只观察。}看清任务做了什么、花了多少、结果是否符合业务要求。一起审阅证据，再决定是否开启运行时控制。\vspace{5pt}}}\par
\vspace{5mm}
\begin{center}
\includegraphics[width=0.72\linewidth]{docs/images/run-analysis.png}\\[2pt]
{\small\color{muted}公开演示：一条多 Agent 任务记录的费用、调用与交接。}
\end{center}
\vspace{0mm}
\sectiontitle{需要准备}
\begin{itemize}
\item Python 3.10+、新的虚拟环境，以及 Reins 私有仓库 \texttt{catyans/reins} 的只读访问权限。
\item 一个真实的多步骤任务、现有测试，以及能定义合格结果的客户工程师。
\item 客户本地的数据目录；模型密钥和原始运行数据保留在客户环境。
\end{itemize}
\sectiontitle{安装试用版}
\texttt{git clone git@github.com:catyans/reins.git /path/to/reins}\\
\texttt{python -m pip install /path/to/reins}

请联系 Reins 团队取得产品私有仓库的访问权限。这份公开资料包提供文档、示例和接入 skill。\par
\vfill
{\small\color{muted}可复制命令和完整说明：\href{https://github.com/catyans/reins-customer-kit}{github.com/catyans/reins-customer-kit}。示例输入和价格均为构造数据。}\par
\newpage

\sectiontitle{1 / 看一次完整任务，不只看模型调用}
在一个完整业务任务外加 \texttt{@trace}，任务开始前只配置一次 Reins。观察模式只记录，不拦截调用。结果是否可用，由原有业务规则判断。

\begin{quote}\ttfamily\small
configure(storage\_path="./pilot.duckdb", mode="observe")\\
@trace(agent\_name="supplier\_agent", task\_type="supplier\_lookup")\\
async def research(source):\\
\hspace*{1em}result = await existing\_agent(source)\\
\hspace*{1em}record\_outcome(success=required\_fields\_pass(result))\\
\hspace*{1em}return result
\end{quote}

对同一输入分别运行接入前后版本，检查 prompt、输出和外部副作用。再执行：

\texttt{reins compare --database ./pilot.duckdb --task-type supplier\_lookup}

设置 \texttt{dashboard=True} 后，在 Agent 运行期间打开本机 \texttt{127.0.0.1:8765} 可看时间线。

\sectiontitle{2 / 接入你想看清的调用}
\renewcommand{\arraystretch}{1.28}
\begin{tabularx}{\linewidth}{@{}p{0.32\linewidth}X@{}}
\toprule
现有调用 & 接入方式 \\
\midrule
OpenAI Chat Completions 或 Anthropic Messages & 任务内的调用可通过受支持的 SDK 路径记录。 \\
原生 Gemini 或付费工具 & 用 \texttt{Workflow.call/acall} 包住每次付费操作，记录真实费用。 \\
框架回调／OTel & 可以补充过程信息；控制支出还需接入实际付费调用。 \\
HTTP Proxy & 可用于实验性的观察。 \\
\bottomrule
\end{tabularx}

\sectiontitle{3 / 再决定要不要加预算规则}
看过真实运行记录后，约定任务预算、每次付费调用的可靠上限、可用模型和工具，以及任务停止时保留什么。先在一小组\textbf{新任务}上试 \texttt{mode="enforce"}。已经发出的调用仍可能计费。

\end{document}
